\documentclass[a4paper,12pt]{article}

% Основные пакеты
\usepackage[utf8]{inputenc}             % Поддержка кодировки UTF-8
\usepackage[T1, T2A]{fontenc}           % Поддержка кириллицы
\usepackage[english, russian]{babel}    % Поддержка русского и английского языков
\usepackage{graphicx}                   % Для вставки изображений
\usepackage{import}                     % Для продвинутого импорта (например, из Inkscape)
\usepackage{subcaption}                 % Для создания подписей к изображениям в сетке
\usepackage{float}                      % Для улучшенного размещения изображений
\usepackage{geometry}                   % Для управления отступами на странице
\usepackage{amsmath, amssymb, amsfonts} % Для математических формул и символов
\usepackage{booktabs}                   % Для красивых таблиц
\usepackage[dvipsnames]{xcolor}         % Для расширенной поддержки цветов
\usepackage[colorlinks]{hyperref}       % Для создания гиперссылок
\usepackage{indentfirst}                % Для красной строки в первом абзаце
\usepackage{listings}                   % Для оформления листингов кода
\usepackage{longtable}                  % Для длинных таблиц
\usepackage{makecell}                   % Для улучшенного форматирования ячеек таблиц
\usepackage{multirow}                   % Для объединения строк в таблицах
\usepackage{tabularray}                 % Для продвинутых таблиц
\usepackage{varwidth}                   % Для переменной ширины блоков
\usepackage{bm}                         % Для жирных математических символов
\usepackage{pifont}                     % Для специальных символов (например, галочек)
\usepackage{fontawesome}                % Для иконок

% Настройка отступов на странице
\geometry{left=2cm, right=2cm, bottom=2cm, top=2cm}

% Настройка гиперссылок
\hypersetup{
	linkcolor=blue,
	filecolor=magenta,
	urlcolor=cyan
}

% Настройка листингов
\definecolor{codegreen}{rgb}{0,0.6,0}
\definecolor{codegray}{rgb}{0.5,0.5,0.5}
\definecolor{codepurple}{rgb}{0.58,0,0.82}
\definecolor{backcolour}{rgb}{0.95,0.95,0.92}

\lstdefinestyle{codestyle}{
	backgroundcolor=\color{backcolour},
	commentstyle=\color{codegreen},
	keywordstyle=\color{blue},
	numberstyle=\tiny\color{codegray},
	stringstyle=\color{magenta},
	basicstyle=\ttfamily\footnotesize,
	breakatwhitespace=false,
	breaklines=true,
	captionpos=b,
	keepspaces=true,
	numbers=left,
	numbersep=5pt,
	showspaces=false,
	showstringspaces=false,
	showtabs=false,
	tabsize=2
}

\lstset{style=codestyle, extendedchars=\true}

\graphicspath{{../images/}}

\begin{document}
	\begin{titlepage}
		\centering
		\vspace*{1cm}
		
		{\large Министерство науки и высшего образования Российской Федерации}\\
		{\large ФЕДЕРАЛЬНОЕ ГОСУДАРСТВЕННОЕ АВТОНОМНОЕ ОБРАЗОВАТЕЛЬНОЕ УЧРЕЖДЕНИЕ ВЫСШЕГО ОБРАЗОВАНИЯ «НАЦИОНАЛЬНЫЙ ИССЛЕДОВАТЕЛЬСКИЙ УНИВЕРСИТЕТ ИТМО»}\\
		{\large (УНИВЕРСИТЕТ ИТМО)}\\
		
		\vspace{2cm}
		
		{\large Факультет «Систем управления и робототехники»}\\
		
		\vspace{3cm}
		
		\textbf{{\Huge ОТЧЕТ}\\
			{\Huge О ЛАБОРАТОРНОЙ РАБОТЕ №5}}\\
		
		\vspace{1cm}
		
		{\LARGE По дисциплине «Техническое зрение»}\\
		{\LARGE на тему: «Преобразование Хафа»}\\
		
		\vspace{2cm}
		
		{\Large Студент:}\\
		Охрименко Ева ИСУ 409290\\
		
		\vspace{2cm}
		
		{\Large Преподаватель:}\\
		Шаветов Сергей Васильевич\\
		
		\vspace{4cm}
		
		{\large г. Санкт-Петербург}\\
		{\large 2025}
		
	\end{titlepage} 
	
	\tableofcontents  % Оглавление
	\newpage
	
	
	
	
\section{Task. Практическое задание}
\subsection{Реализация методов для численной оценки градиента и гессиана}
\subsubsection{Градиент}

Формула для каждой компоненты градиента:

\[
\frac{\partial f}{\partial x_i} \approx \frac{f(x_1,\dots,x_i+h,\dots,x_n) - f(x_1,\dots,x_n)}{h}
\]

Где:
\begin{itemize}
	\item $h$ - эпсилон
	\item $f$ - наша функция
	\item $x_i$ - текущая координата
\end{itemize}

Весь градиент это вектор таких производных:

\[
\nabla f \approx \left(
\frac{f(x+h,y)-f(x,y)}{h}, 
\frac{f(x,y+h)-f(x,y)}{h}
\right)
\]
Вот код, который реализует это:

\begin{lstlisting}[language=python, caption={Градиент}]
def numericalGradient(f, x, h=1e-5):
	grad = np.zeros_like(x)
	fx = f(x)
	for i in range(len(x)):
		xT = x.copy()
		xT[i] += h
		grad[i] = (f(xT) - fx) / h
	return grad
\end{lstlisting}



\subsubsection{Гессиан}

Мы хотим найти матрицу вторых производных (гессиан) функции $f(x_1, x_2, ..., x_n)$ в точке.

Для этого мы:
\begin{itemize}
	\item Берем маленький шаг $h$ - эпсилон
	\item Смотрим, как меняется функция при небольших изменениях по разным направлениям
	\item Используем комбинации значений функции в соседних точках
\end{itemize}

Мы считаем так:

Для каждой пары переменных $(x_i, x_j)$:

1. Создаем 4 точки, немного меняя $x_i$ и $x_j$:
\begin{itemize}
	\item Точка 1: $x_i + h$, $x_j + h$
	\item Точка 2: $x_i + h$, $x_j - h$
	\item Точка 3: $x_i - h$, $x_j + h$
	\item Точка 4: $x_i - h$, $x_j - h$
\end{itemize}

2. Считаем элемент матрицы по формуле:
\[
H_{ij} = \frac{f(\text{Точка1}) - f(\text{Точка2}) - f(\text{Точка3}) + f(\text{Точка4})}{4h^2}
\]

3. Так как матрица симметричная, $H_{ji} = H_{ij}$

Для функции $f(x,y)$ гессиан будет:
\[
H = \begin{pmatrix}
	\frac{\partial^2 f}{\partial x^2} & \frac{\partial^2 f}{\partial x \partial y} \\
	\frac{\partial^2 f}{\partial y \partial x} & \frac{\partial^2 f}{\partial y^2}
\end{pmatrix}
\]

Вот код:
\begin{lstlisting}[language=python, caption={Гессиан}]
def numericalHessian(f, x, h=1e-5):
	n = len(x)
	hess = np.zeros((n, n))
	for i in range(n):
		for j in range(i, n):x1, x2 = x.copy(), x.copy()
		x1[i] += h; x1[j] += h
		x2[i] += h; x2[j] -= h
		x3, x4 = x.copy(), x.copy()
		x3[i] -= h; x3[j] += h
		x4[i] -= h; x4[j] -= h
		hess[i,j] = (f(x1) - f(x2) - f(x3) + f(x4)) / (4 * h**2)
		hess[j,i] = hess[i,j]
			
	return hess
\end{lstlisting}

\subsection{Оценка зависимости точности градиента и гессиана от выбора шага h}



\end{document}